% The causal response as a gap between two paths in the hand economy. The
% counterfactual switches off the t = 5 intervention and keeps every disturbance.
% Blue brackets are realized minus counterfactual (1, 0.5, 0.25); grey arrows
% are the change since t = 4, which mixes the intervention with everything else.
\documentclass[border=3pt]{standalone}
% ————— Shared preamble for the course figures —————
% Each figure is a standalone TikZ document. build.sh compiles it to a PDF (for
% the LaTeX build) and converts that to an SVG (for the web build).
%
% Nothing is drawn in pure black: the chrome sits at a mid grey that stays
% legible on white paper and on the site's dark theme alike. Data colours are
% the course palette, the same blue and orange used in the practicum.

\usepackage{amsmath}
\usepackage{tikz}
\usepackage{pgfplots}
\pgfplotsset{compat=1.18}
\usepgfplotslibrary{fillbetween,groupplots}
\usetikzlibrary{arrows.meta,positioning,calc,backgrounds}

\definecolor{nwBlue}{HTML}{2A78D6}
\definecolor{nwOrange}{HTML}{D9601F}
\definecolor{nwMut}{HTML}{898781}
\definecolor{nwInk}{HTML}{6B6965}

\pgfplotsset{
  nwaxis/.style={
    axis lines=left,
    axis line style={nwMut, line width=0.5pt},
    tick align=outside,
    tick style={nwMut, line width=0.5pt},
    label style={text=nwInk, font=\small},
    tick label style={text=nwMut, font=\footnotesize},
    legend cell align=left,
    legend style={draw=none, fill=none, font=\footnotesize, text=nwInk,
                  row sep=1pt, inner sep=1pt},
    every axis plot post/.append style={line join=round},
    clip mode=individual,
  },
}

\directlua{dofile("fig-l01-realized-counterfactual-data.lua")}
\begin{document}
\begin{tikzpicture}[
  gap/.style={nwBlue, line width=1.1pt, {Bar[width=4.5pt]}-{Bar[width=4.5pt]}},
  rise/.style={nwMut, line width=0.7pt, -{Stealth[length=4.5pt,width=3.5pt]}},
]
\begin{axis}[
  nwaxis,
  width=12cm, height=7cm,
  xmin=0.5, xmax=12.5, ymin=-2.25, ymax=2.95,
  xtick={1,2,...,12},
  ytick={-2,-1,0,1,2},
  xlabel={calendar time $t$},
  ylabel={outcome},
]

% the level at t = 4, from which a "change since" is measured
\addplot[nwMut, densely dotted, line width=0.7pt]
  coordinates {(4,\lpRcYd) (7.5,\lpRcYd)};
\node[text=nwMut, font=\footnotesize, anchor=north east]
  at (axis cs:3.92,\lpRcYd-0.04) {$y_4$};

% the two paths
\addplot[nwOrange, line width=1.1pt, dash pattern=on 4pt off 2.5pt,
         mark=o, mark size=1.8pt, mark options={solid, fill=white}]
  coordinates {\lpRcYc};
\addplot[nwInk, line width=1.1pt, mark=*, mark size=1.7pt,
         mark options={fill=nwInk, draw=nwInk}]
  coordinates {\lpRcY};

% realized minus counterfactual
\draw[gap] (axis cs:4.78,\lpRcYce) -- (axis cs:4.78,\lpRcYe);
\draw[gap] (axis cs:6.22,\lpRcYcf) -- (axis cs:6.22,\lpRcYf);
\draw[gap] (axis cs:7.22,\lpRcYcg) -- (axis cs:7.22,\lpRcYg);
\node[text=nwBlue, font=\footnotesize, anchor=east]
  at (axis cs:4.68,{0.5*(\lpRcYce+\lpRcYe)}) {$1$};
\node[text=nwBlue, font=\footnotesize, anchor=west]
  at (axis cs:6.32,{0.5*(\lpRcYcf+\lpRcYf)}) {$0.5$};
\node[text=nwBlue, font=\footnotesize, anchor=west]
  at (axis cs:7.32,{0.5*(\lpRcYcg+\lpRcYg)}) {$0.25$};

% change since t = 4
\draw[rise] (axis cs:5.2,\lpRcYd) -- (axis cs:5.2,\lpRcYe);
\draw[rise] (axis cs:6.8,\lpRcYd) -- (axis cs:6.8,\lpRcYg);
\node[text=nwMut, font=\footnotesize, anchor=south east]
  at (axis cs:5.14,\lpRcYd+0.06) {$-0.731$};
\node[text=nwMut, font=\footnotesize, anchor=south east]
  at (axis cs:6.92,\lpRcYg+0.06) {$+1.970$};

% direct labels for the paths
\node[text=nwInk, font=\footnotesize, anchor=south west]
  at (axis cs:1.45,1.98) {realized $y_t$};
\node[text=nwOrange, font=\footnotesize, anchor=north west]
  at (axis cs:4.35,-1.8) {counterfactual $y^{c}_t$: the $t=5$ intervention switched off};

% key for the two kinds of vertical mark
\draw[gap]  (axis cs:0.95,-1.33) -- (axis cs:0.95,-1.03);
\node[text=nwBlue, font=\footnotesize, anchor=west]
  at (axis cs:1.12,-1.18) {$y_t-y^{c}_t$};
\draw[rise] (axis cs:0.95,-1.53) -- (axis cs:0.95,-1.83);
\node[text=nwMut, font=\footnotesize, anchor=west]
  at (axis cs:1.12,-1.68) {$y_t-y_4$};

\end{axis}
\end{tikzpicture}
\end{document}
